\chapter{Training CSC-LISTA}\label{ch:training}

\section{One training step}

\begin{enumerate}
\item Pick 8 random 512 pixel crops of simulated fields. A crop must contain real signal: its peak
must be at least 0.05 in normalised units, or the code tries another position.
\item Pad each input to 768 pixels.
\item Run the network.
\item Compute the loss on the middle of the crop. A margin of 64 pixels is dropped on each side,
because edges are the least reliable.
\item Backpropagate, clip the gradient, take an Adam step.
\end{enumerate}

\begin{code*}{train_csc_lista.py}
            y = model(b)[..., pad:-pad, pad:-pad]
            keep = (slice(None), slice(margin, -margin), slice(margin, -margin))
            loss = crop_loss(y[keep], target[keep], args.loss)
        optimiser.zero_grad()
        loss.backward()
        torch.nn.utils.clip_grad_norm_(model.parameters(), 1.0)
        optimiser.step()
        schedule.step()
\end{code*}

\section{The loss: MSE or relative}

\begin{code*}{train_csc_lista.py}
def crop_loss(output, target, kind):
    """Mean squared error, or the mean per-crop relative L2 used for validation."""
    if kind == 'relative':
        return relative_l2(output, target).mean()
    return ((output - target) ** 2).mean()
\end{code*}

\textbf{MSE} is the paper's loss. It counts every pixel's squared error equally, so crops with strong fields
dominate. \textbf{Relative} divides by each crop's own size, so all crops count equally. It is also exactly the
metric we report. We tried it because of a problem you will see in Chapter~\ref{ch:density}.

\section{The learning rate problem}

The first attempt crashed with NaN. The filters start around $2\times10^{-8}$ in size, and a normal Adam
rate of $10^{-4}$ means each step is thousands of times larger than the value it changes. So the rate is now set
relative to each parameter's own scale.

\begin{code*}{train_csc_lista.py}
    # Filters start near 1e-8, so a fixed rate would swamp them; scale it to each.
    groups = [dict(params=[model.w], lr=args.lr_relative * float(model.w.pow(2).mean().sqrt())),
              dict(params=[model.theta], lr=args.lr_relative * float(model.theta.abs().mean()))]
\end{code*}

With \code{lr_relative} at 0.03, each step changes a parameter by about 3\% of the typical size. The rate then
decays along a cosine curve to zero by the last step.

\section{Validation and checkpoints}

Every 100 or 250 steps, the model is scored on 64 fixed crops from the validation layouts, which are
never trained on. If the score beats the best so far, the model is saved to \code{best.pt}. After training,
that best checkpoint is loaded and used for the COMSOL score.

\begin{code*}{train_csc_lista.py}
            if step == 0 or score < best:
                best = score
                torch.save(dict(state=model.state_dict(), step=step, score=score),
                           output / 'best.pt')
\end{code*}

This matters a lot. In the strict runs the best step came very early, at 300 of 1,000 and 200 of 1,000.

\begin{careful}[Validation and COMSOL scores cannot be compared directly]
Validation crops are chosen to contain signal, and COMSOL scores use the full image, including large
blank areas. So ``48.8\% on validation and 33.8\% on COMSOL'' does not mean COMSOL is easier.
\end{careful}

\section{Does the filter size matter?}

You asked whether to make $w$ bigger. The physics suggested it might, since the blurred kernel has a long
tail. Three 800-step trials on simulated validation gave a clear answer.

\begin{center}
\begin{tabular}{cc}
\toprule
Filter width & Validation error \\
\midrule
155 px & 52.35\% \\
255 px & 52.02\% \\
355 px & 51.99\% \\
\bottomrule
\end{tabular}
\end{center}

The difference is about a third of a point, and larger filters cost much more. So 155 stayed.

\section{Hardware}

All runs used the RTX A6000 with 48 GB, at batch 8.

\begin{center}
\begin{tabular}{rrrr}
\toprule
Layers & Parameters & Peak GPU memory & Time \\
\midrule
16 & 12.3 million & 12.8 GiB & 10,940 s for 2,000 steps \\
32 & 24.6 million & 22.0 GiB & 5,037 s for 1,000 steps \\
64 & 49.2 million & 40.2 GiB & 10,058 s for 1,000 steps \\
\bottomrule
\end{tabular}
\end{center}

Memory and time both grow with the number of layers, as you would expect from a chain of identical
layers. Also note the two lengths: the first runs took 2,000 steps. Once MSE was seen to peak early, later
runs used 1,000.

\begin{careful}[One seed]
Every run here used a single random seed. A difference of less than a point or so between two runs may
be luck. No error bars exist.
\end{careful}
